\section{Related Work}
\label{sec:related}

\paragraph{Sound predictive analysis for shared memory.} Predictive analysis takes one
trace of a multithreaded program and proves that a reordering of it is a run of every
program consistent with the trace. WCP~\citep{kini17wcp} gave the first sound partial
order weaker than happens-before with a linear-time algorithm. SHB~\citep{mathur18shb}
showed that happens-before is sound only up to the first race and repaired it.
SDP~\citep{genc19sdp}, M2~\citep{pavlogiannis20m2} and sync-preserving
races~\citep{mathur21syncp} weakened the order further. SyncP is sound and complete for
the reorderings that preserve the order of critical sections, which is the template for
\cref{thm:complete}. Deadlock prediction followed the same path, from executability
constraints solved by SMT~\citep{kalhauge18sound} to sync-preserving deadlocks in linear
time~\citep{tunc23spd}, the template for \cref{thm:l1dead}. Synchronisation-reversal
races~\citep{shi24osr} predict races that need a lock order reversed; reversing one match
is the message-passing analogue of that step. Strong trace prefixes~\citep{ang24strong}
let a monitor drop parts of the trace that cannot affect the property. Our horizon is a
strong prefix for the matching alphabet, with program order, deliveries, collectives and
rendezvous as strong dependencies and no weak ones. A unified definition of soundness for
these analyses is due to~\citet{liu25soundness}. It is stated for shared memory with locks,
and we restate it for matching consistency. Tree clocks~\citep{mathur22treeclock} speed up
the causal orders these analyses compute. The complexity of race
prediction~\citep{mathur20lics} and its fine-grained lower bounds~\citep{kulkarni21concur}
frame the open question of choosing which receives to change.

\paragraph{Complexity of message-passing consistency.} \citet{shi26popl} map the
complexity of deciding whether an abstract execution of a channel-based program is
consistent, with or without a fixed reads-from relation. Our check is their VCh-rf problem
for MPI's channel shape. Queues are single-producer and single-consumer per (source,
destination, communicator), wildcard receives take from a union of queues, buffering is a
per-message threshold, and collectives synchronise by shape. Their linear case requires
synchronous channels. MPI's structure keeps the problem tractable for mixed eager and
rendezvous traffic. They also use swaps of two reads-from pairs as a mutation operator,
which we use as one of two alternative shapes.

\paragraph{Prediction for MPI.} MOPPER~\citep{forejt14mopper,forejt17mopper} proved
deadlock detection NP-complete for single-path MPI programs and encoded it
propositionally. Later work from the same circle encodes candidates for an SMT
solver~\citep{huang20ase} and compresses traces to scale it~\citep{huang23tpds}. Another
group refines the match sets~\citep{cluster22matchsets}, partitions programs with logical
clocks~\citep{li24parco}, and compresses paths and focuses the
matching~\citep{internetware24}. All of them target deadlocks under the
single-path assumption, which \cref{lem:sp} shows unsound on metadata-dependent code.
None states a soundness theorem for the alternative itself or a buffering profile. The
idea of reasoning about the part of a message-passing execution unaffected by an earlier
race is older. The MAD environment stated it for message races~\citep{kranzlmueller00mad},
and the unaffected-race papers define races causally and detect them on the
fly~\citep{park04unaffected,park07mpirace}. These works do not certify that an
alternative is executable.

\paragraph{Exhaustive exploration.} ISP~\citep{vakkalanka08isp}, DAMPI~\citep{vo10dampi},
MPI-SV~\citep{yu20mpisv} and CIVL~\citep{siegel15civl} explore the schedules of a program
or verify it symbolically. They see everything a program can do on an input, but they pay
for every competing wildcard. We use re-execution only to confirm. The stress-testing
baseline of \cref{sec:eval} follows the randomised concurrency testing
line~\citep{wolff24rff,zhao25surw}. State-space estimation for DPOR~\citep{balasubramanian26dpor}
quantifies how much exploration a program needs.

\paragraph{Observational tools and replay.} \MUST~\citep{hilbrich12must} checks the
recorded run for argument errors, type mismatches, deadlocks and, with ThreadSanitizer,
buffer races. PARCOACH~\citep{saillard14parcoach} checks collective ordering statically
and dynamically. Record and replay for MPI~\citep{netzer92replay,sato15rempi} reproduces
the matching of a recorded run. Our pins force a matching the run did \emph{not} take.

\paragraph{Message passing beyond MPI.} Vector-clock analyses of Go channels report
alternative communications~\citep{sulzmann18go}, and dl$^2$ finds communication deadlocks
in distributed deep-learning jobs observationally~\citep{gao25dl2}. MPI's matching
semantics (envelopes, wildcards, non-overtaking, buffering regimes, collectives)
specialises to Go's channels, not the other way round.

\paragraph{Benchmarks.} The MPI Bugs Initiative~\citep{laurent21mbi} and
\MBB~\citep{jammer24mbb} are the standard corpora. \MBB merges earlier suites and reports
that \MUST is the strongest existing tool.
